%% This is a LaTeX document.  Hey, Emacs, -*- latex -*- , get it?
\newif\ifbeamermode
\beamermodetrue
\ifbeamermode
\documentclass[mathserif,a4paper,aspectratio=169]{beamer}
\else
\documentclass[a4paper]{article}
\usepackage{beamerarticle}
\usepackage[a4paper,margin=2.5cm]{geometry}
\fi
\usepackage[shorthands=off,french]{babel}
\usepackage[utf8]{inputenc}
\usepackage[T1]{fontenc}
\usepackage{lmodern}
\DeclareUnicodeCharacter{00A0}{~}
\DeclareUnicodeCharacter{2026}{...}
% Beamer theme:
\usetheme{Goettingen}
%\usecolortheme{albatross}
%\usecolortheme{lily}
%\setbeamercovered{transparent}
% A tribute to the worthy AMS:
\usepackage{amsmath}
\usepackage{amsfonts}
\usepackage{amssymb}
\usepackage{amsthm}
%
\usepackage{mathrsfs}
\usepackage{mathpartir}
\usepackage{flagderiv}
%
\usepackage{graphicx}
\usepackage{tikz}
\usetikzlibrary{arrows,automata,calc}
\usepackage{hyperref}
%
\newcommand{\itempoint}{\strut\hbox{\color{beamerstructure}\donotcoloroutermaths$\blacktriangleright$}\nobreak\hskip.5em plus.5em\relax}
\renewcommand{\thefootnote}{\textdagger}
\newcommand{\dbllangle}{\mathopen{\langle\!\langle}}
\newcommand{\dblrangle}{\mathclose{\rangle\!\rangle}}
\newcommand{\cps}{\mathrm{CPS}}
%\newcommand{\dottedlimp}{\mathbin{\dot\Rightarrow}}
%\newcommand{\dottedland}{\mathbin{\dot\land}}
%\newcommand{\dottedlor}{\mathbin{\dot\lor}}
%\newcommand{\dottedtop}{\mathord{\dot\top}}
%\newcommand{\dottedbot}{\mathord{\dot\bot}}
%\newcommand{\dottedneg}{\mathop{\dot\neg}}
\mathchardef\emdash="07C\relax
\newcommand{\mpdotsabove}[1]{\inferrule*{\vdots}{#1}}
%\setlength{\derivskip}{4pt}
%
\mode<article>{
    \newcommand{\donotcoloroutermaths}{\relax}
    \renewcommand{\frametitle}[1]{\subsubsection*{#1 {\footnotesize[transp. \insertframenumber]}}}
    \newcommand{\myitemmark}{\hbox{$\blacktriangleright$}}
    \renewcommand{\itempoint}{\strut\myitemmark\nobreak\hskip.5em plus.5em\relax}
    \setlength{\parindent}{0pt}
    \renewcommand{\thepage}{III–\arabic{page}}
}
%
%
%
\title{{\footnotesize Partie III:} Introduction aux quantificateurs}
\subtitle{CSC-3TC34-TP / INF110 (Logique et Fondements de l'Informatique)}
\author[David Madore]{David A. Madore\\
{\footnotesize Télécom Paris}\\
\texttt{david.madore@}(nom de l'école)\texttt{.fr}}
\date{2023–2026}
\mode<presentation>{%
\beamertemplatenavigationsymbolsempty
\usenavigationsymbolstemplate{\vbox{\hbox{\footnotesize(III) \hyperlinkslideprev{$\leftarrow$}\insertframenumber/\inserttotalframenumber\hyperlinkslidenext{$\rightarrow$}}}}
}
\setbeamercolor{myhighlight}{fg=black,bg=white!90!green}
\colorlet{mydarkgreen}{green!50!black}
\begin{document}
\mode<article>{\maketitle\renewcommand{\labelitemi}{\myitemmark}}
%
\setlength\abovedisplayskip{2pt plus 2pt minus 2pt}
\setlength\belowdisplayskip{2pt plus 2pt minus 2pt}
%
\begin{frame}
\titlepage
{\footnotesize\center{\url{http://perso.enst.fr/madore/inf110/transp-inf110.pdf}}\par}
{\tiny
\immediate\write18{sh ./vc > vcline.tex}
\begin{center}
Git: \input{vcline.tex}
\end{center}
\immediate\write18{echo ' (stale)' >> vcline.tex}
\par}
\end{frame}
%
\section*{Plan}
\begin{frame}
\frametitle{Plan}
\tableofcontents
\end{frame}
%
\section{Les quantificateurs : discussion informelle}
\begin{frame}
\frametitle{Limitations du calcul propositionnel}

\itempoint On a parlé pour l'instant de \textbf{calcul
  propositionnel}, qui ne connaît que les affirmations logiques et les
connecteurs propositionnels $\Rightarrow,\land,\lor,\top,\bot$.

\medskip

\itempoint Mais il y a deux notations logiques essentielles en
mathématiques au-delà de ces connecteurs : les
\textbf{quantificateurs} $\forall,\exists$, qui :
\begin{itemize}
\item prennent une formule $P(v)$ dépendant d'une variable $v$ libre
  {\footnotesize (de type $I$)},
\item lient cette variable, pour former une nouvelle formule :
  $\forall (v:I). P(v)$ ou $\exists (v:I). P(v)$
  {\footnotesize (parfois juste $\forall v. P(v)$ et $\exists v. P(v)$)}.
\end{itemize}

\medskip

\itempoint Intuitivement, il faut penser à $\forall$ et $\exists$
comme des « $\land$ et $\lor$ en famille », c'est-à-dire que :
\begin{itemize}
\item $\forall v.P(v)$, parfois noté $\bigwedge_v P(v)$ est à $P\land
  Q$ ce que $\prod_i p_i$ est à $p\times q$,
\item $\exists v.P(v)$, parfois noté $\bigvee_v P(v)$ est à $P\lor Q$
  ce que $\sum_i p_i$ est à $p + q$.
\end{itemize}

\medskip

\itempoint Il existe de \alert{nombreux systèmes logiques} différant
notamment en \alert{ce qu'on a le droit de quantifier} (qui sont les
$v$ ici ? quel est leur domaine $I$ ?).

\smallskip

\textcolor{brown}{Commençons par une discussion informelle de
  $\forall$ et $\exists$.}

\end{frame}
%
\begin{frame}
\frametitle{L'interprétation BHK des quantificateurs}

On a déjà vu l'interprétation informelle des connecteurs, on introduit
maintenant les quantificateurs :

\begin{itemize}
\item {\color{darkgray} un témoignage de $P\land Q$, est un témoignage
  de $P$ et un de $Q$,}
\item {\color{darkgray} un témoignage de $P\lor Q$, est un témoignage
  de $P$ ou un de $Q$, et la donnée duquel des deux on a choisi,}
\item {\color{darkgray} un témoignage de $P\Rightarrow Q$ est un moyen
  de transformer un témoignage de $P$ en un témoignage de $Q$,}
\item {\color{darkgray} un témoignage de $\top$ est trivial,} \quad
  \itempoint {\color{darkgray} un témoignage de $\bot$ n'existe pas,}
\item un témoignage de $\forall v.P(v)$ est un moyen
  de transformer un $x$ quelconque en un témoignage de $P(x)$,
\item un témoignage de $\exists v.P(v)$ est la donnée d'un certain
  $t_0$ et d'un témoignage de $P(t_0)$.
\end{itemize}

\end{frame}
%
\begin{frame}
\frametitle{Curry-Howard pour le $\forall$}

\itempoint On a vu que Curry-Howard fait correspondre
\alert{conjonction logique} $P\land Q$ {\footnotesize (« un témoignage
  de $P$ et un de $Q$ »)} avec \alert{type produit} $\sigma\times\tau$
      {\footnotesize (« une valeur de $\sigma$ et une de $\tau$ »)}.

\medskip

\itempoint De façon analogue, la \alert{quantification universelle}
$\forall v. P(v)$ {\footnotesize (« une façon de transformer $v$ en un
  témoignage de $P(v)$ »)}, qui est une sorte de \emph{conjonction en
famille} $\bigwedge_v P(v)$, correspondra au \alert{type produit en
  famille} $\prod_v \sigma(v)$ {\footnotesize (« fonction renvoyant
  pour chaque $v$ une valeur de $\sigma(v)$ »)}.

\medskip

\itempoint Ceci présuppose l'existence de \alert{familles de types} $v
\mapsto \sigma(v)$ (= types dépendant de quelque chose) dont on puisse
prendre le produit.

\medskip

\itempoint Une preuve de $\forall (v:I).P(v)$ correspondra à un terme
de forme $\lambda(v:I).\,(\cdots)$, où le type de $(\cdots)$
correspond à $P(v)$.

\medskip

\itempoint Remarquer que $\forall (v:I).P$, si $P$ ne dépend pas de
$v$, « ressemble » à $I \Rightarrow P$ de la même manière que
$\prod_{i\in I} S = S^I$ (ensemblistement ou numériquement).
           {\footnotesize (Les détails dépendent de la nature de la
             quantification.)}

\end{frame}
%
\begin{frame}
\frametitle{Curry-Howard pour le $\exists$}

\itempoint On a vu que Curry-Howard fait correspondre
\alert{disjonction logique} $P\lor Q$ {\footnotesize (« un témoignage
  de $P$ ou un de $Q$, avec la donnée duquel on a choisi »)} avec
\alert{type somme} $\sigma+\tau$ {\footnotesize (« une valeur de
  $\sigma$ ou une de $\tau$, avec un sélecteur entre les deux »)}.

\medskip

\itempoint De façon analogue, la \alert{quantification existentielle}
$\exists v. P(v)$ {\footnotesize (« la donnée d'un $t_0$ et d'un
  témoignage de $P(t_0)$ »)}, qui est une sorte de \emph{disjonction
en famille} $\bigvee_v P(v)$, correspondra au \alert{type somme en
  famille} $\sum_v \sigma(v)$ {\footnotesize (« donnée d'un $t_0$ et
  d'une valeur de type $\sigma(t_0)$ »)}.

\medskip

\itempoint Une preuve de $\exists (v:I).P(v)$ correspondra à un terme de
forme $\langle t_0, \cdots\rangle$, où le type de $(\cdots)$
correspond à $P(t_0)$.  {\footnotesize (De nouveau, il faut des
  « familles de types ».)}

\medskip

\itempoint Remarquer que $\exists (v:I).P$, si $P$ ne dépend pas de $v$,
« ressemble » à $I \times P$ de la même manière que $\sum_{i\in I} S =
I\times S$.  {\footnotesize (Les détails dépendent de la nature de la
  quantification.)}

\medskip

\itempoint Mais Curry-Howard atteint ses limites : il n'est pas dit
que d'une preuve de $\exists v.P(v)$ on \alert{puisse extraire} le
$t_0$ correspondant dans autre chose qu'une preuve.  {\footnotesize
  (Les détails dépendent du système logique précis considéré {\tiny et
    si Martin-Löf est dans la salle}.)}

\end{frame}
%
\begin{frame}
\frametitle{Règles d'introduction et d'élimination de $\forall$}

{\footnotesize\textcolor{orange}{Les règles ci-dessous (et
    transp. suivants) sont \alert{incomplètes} : il manque des
    explications sur le type $I$ sur lequel on quantifie et comment on
    peut en former des « termes d'individus ».}\par}

\medskip

\itempoint \underline{Introduction du $\forall$ :} pour montrer
$\forall (v:I).\, Q$, on s'arrange (quitte à renommer la variable
liée) pour que $v:I$ soit « frais », c'est-à-dire qu'il n'apparaisse
(libre) dans \alert{aucune hypothèse} en cours : si on montre $Q$ sur
ce $v$ « arbitraire », on peut conclure $\forall (v:I).\, Q$.

\smallskip

{\footnotesize (Rédaction : « soit $v$ arbitraire (…) on a $Q(v)$ ;
  donc $\forall (v:I).\, Q(v)$ ».)\par}

\medskip

\textcolor{teal}{Ceci donnera un $\lambda$-terme noté
  $\lambda(v:I).(\cdots)$ comme l'ouverture d'une hypothèse.}

\bigskip

\itempoint \underline{Élimination du $\forall$ :} pour utiliser
$\forall (v:I). Q$, on peut l'appliquer à un $t$ quelconque (un
\alert{terme} de type $I$).

\smallskip

{\footnotesize (Rédaction : « on a $\forall (v:I). Q(v)$ et $t$ de
  type $I$ ; en particulier, on a $Q(t)$ ».)\par}

\medskip

\textcolor{teal}{Ceci donnera un $\lambda$-terme noté $ft$ comme
  l'application d'une implication.}

\end{frame}
%
\begin{frame}
\frametitle{Règles d'introduction et d'élimination de $\exists$}

\itempoint \underline{Introduction du $\exists$ :} pour montrer
$\exists (v:I). Q$, on peut le montrer sur un terme $t$ quelconque de
type $I$.

\smallskip

{\footnotesize (Rédaction : « pour ce $t$ de type $I$ on a $Q(t)$ ; en
  particulier, on a $\exists (v:I). Q(v)$ ».)\par}

\medskip

\textcolor{teal}{Ceci donnera un $\lambda$-terme noté $\langle
  t,\cdots\rangle$ comme pour une conjonction.}

\bigskip

\itempoint \underline{Élimination du $\exists$ :} pour utiliser
$\exists (v:I). P(v)$ pour montrer une conclusion $Q$, on s'arrange
(quitte à renommer la variable liée) pour que $v$ soit « frais »,
c'est-à-dire qu'il n'apparaisse (libre) dans \alert{aucune hypothèse}
en cours \alert{ni dans la conclusion} $Q$ : si on montre $Q$ à partir
de $P$ sur ce $v$ « arbitraire », on peut conclure $Q$ à partir de
$\exists v. P$.

\smallskip

{\footnotesize (Rédaction : « on a $\exists (v:I). P(v)$ : soit $v$
  arbitraire tel que $P(v)$ : (…) on a $Q$ ; donc $Q$ ».)\par}

\medskip

\textcolor{teal}{Ceci donnera un $\lambda$-terme noté
  $(\texttt{match~}\cdots\texttt{~with~}\langle v,h\rangle \mapsto
  \cdots)$.}

\end{frame}
%
\begin{frame}
\label{natural-deduction-rules-with-quantifiers}
\frametitle{Aperçu d'ensemble des règles de la déduction naturelle}

\begin{tabular}{c|c|c}
&Intro&Élim\\\hline
$\Rightarrow$
&\scalebox{0.65}{$\inferrule{\Gamma,{\color{mydarkgreen}P}\vdash Q}{\Gamma\vdash P\Rightarrow Q}$}
&\scalebox{0.65}{$\inferrule{\Gamma\vdash P\Rightarrow Q\\\Gamma\vdash P}{\Gamma\vdash Q}$}
\\\hline
$\land$
&\scalebox{0.65}{$\inferrule{\Gamma\vdash Q_1\\\Gamma\vdash Q_2}{\Gamma\vdash Q_1\land Q_2}$}
&\scalebox{0.65}{$\inferrule{\Gamma\vdash Q_1\land Q_2}{\Gamma\vdash Q_1}$}
\quad\scalebox{0.65}{$\inferrule{\Gamma\vdash Q_1\land Q_2}{\Gamma\vdash Q_2}$}
\\\hline
$\lor$
&\scalebox{0.65}{$\inferrule{\Gamma\vdash Q_1}{\Gamma\vdash Q_1\lor Q_2}$}
\quad\scalebox{0.65}{$\inferrule{\Gamma\vdash Q_2}{\Gamma\vdash Q_1\lor Q_2}$}
&\scalebox{0.65}{$\inferrule{\Gamma\vdash P_1\lor P_2\\\Gamma,{\color{mydarkgreen}P_1}\vdash Q\\\Gamma,{\color{mydarkgreen}P_2}\vdash Q}{\Gamma\vdash Q}$}
\\\hline
$\top$
&\scalebox{0.65}{$\inferrule{\strut}{\Gamma\vdash \top}$}
&\scalebox{0.65}{(néant)}
\\\hline
$\bot$
&\scalebox{0.65}{(néant)}
&\scalebox{0.65}{$\inferrule{\Gamma\vdash \bot}{\Gamma\vdash Q}$
(ou pour la logique classique : $\inferrule{\Gamma,\neg Q\vdash \bot}{\Gamma\vdash Q}$)}
\\\hline
$\forall$
&$\inferrule{\Gamma, {\color{mydarkgreen}v:I}\vdash Q}{\Gamma\vdash \forall (v:I).\, Q}$ ($v$ \alert{frais})
&$\inferrule{\Gamma\vdash \forall (v:I).\, Q\\{\color{gray}\Gamma\vdash t:I}}{\Gamma\vdash Q[v\backslash t]}$
\\\hline
$\exists$
&$\inferrule{{\color{gray}\Gamma\vdash t:I}\\\Gamma\vdash Q[v\backslash t]}{\Gamma\vdash \exists (v:I).\, Q}$
&$\inferrule{\Gamma\vdash \exists (v:I).\, P\\\Gamma, {\color{mydarkgreen}v:I}, {\color{mydarkgreen}P}\vdash Q}{\Gamma\vdash Q}$
($v$ \alert{frais})
\\
\end{tabular}

\smallskip

\itempoint « $v$ frais » = « $v$ n'apparaît nulle part ailleurs »
(cf. transp. précédents).

\smallskip

\itempoint Le contexte $\Gamma$ peut contenir des formules
(hypothèses) et des variables d'« individus » (avec leur type,
p.ex. $v:I$).

\end{frame}
%
\begin{frame}
\frametitle{Notations des $\lambda$-termes pour $\forall$}

{\footnotesize

\itempoint On a vu en calcul propositionnel qu'on peut noter les
démonstrations par des « $\lambda$-termes » qui peuvent ensuite être
réinterprétés comme des programmes (c'est Curry-Howard).  Complétons
ces notations pour $\forall, \exists$ :

\par}

\medskip

\itempoint \textcolor{blue}{\underline{Introduction du $\forall$ :}}
si $s$ désigne une preuve de $Q$ faisant intervenir $v$ variable libre
de type $I$, on notera $\lambda(v:I).\,s$ la preuve de
$\forall(v:I).\,Q$ obtenue par introduction du $\forall$.
{\footnotesize (\textbf{N.B.} $v$ peut apparaître dans $Q$ mais pas
  dans $\Gamma$.)}
\[
\inferrule{\Gamma, v:I\vdash s:Q}{\Gamma\vdash \lambda(v:I).\,s\;:\;\forall (v:I).\, Q}
\]

\medskip

\itempoint \textcolor{blue}{\underline{Élimination du $\forall$ :}} si
$f$ désigne une preuve de $\forall(v:I).\,Q$ et $t$ un terme de
type $I$, on notera $ft$ la preuve de $Q[v\backslash t]$ {\footnotesize
  (c'est-à-dire $Q$ avec $v$ remplacé par $t$)} obtenue par
élimination du $\forall$ sur ce terme.

{\footnotesize (\textbf{N.B.} on n'explique pas comment le « terme
  d'individu » $t$ est formé.)}
\[
\inferrule{\Gamma\vdash f \;:\; \forall (v:I).\, Q\\{\color{gray}\Gamma\vdash t:I}}{\Gamma\vdash ft \;:\; Q[v\backslash t]}
\]

\medskip

\itempoint Ceci est conforme à l'idée de BHK : une preuve de
$\forall(v:I).\, Q(v)$ prend un $x$ de type $I$ et renvoie une preuve
de $Q(x)$.

\end{frame}
%
\begin{frame}
\frametitle{Notations des $\lambda$-termes pour $\exists$}

\itempoint \textcolor{blue}{\underline{Introduction du $\exists$ :}}
si $t$ désigne un terme de type type $I$ et $z$ une preuve de
$Q[v\backslash t]$ (pour ce $t$-là, donc), on notera $\langle
t,z\rangle$ la preuve de $\exists(v:I).\,Q$ obtenue par introduction
du $\exists$.
\[
\inferrule{{\color{gray}\Gamma\vdash t:I}\\\Gamma\vdash z \;:\; Q[v\backslash t]}{\Gamma\vdash \langle t,z\rangle \;:\; \exists (v:I).\, Q}
\]

\medskip

\itempoint \textcolor{blue}{\underline{Élimination du $\exists$ :}} si
$z$ désigne une preuve de $\exists(v:I).\,P$ et $s$ une preuve,
faisant intervenir $v$ variable libre de type $I$, de $Q$ qui ne fait
pas intervenir $v$, et $h$ hypothèse supposant $P$ (pour ce $v$-là,
donc), on notera $(\texttt{match~}z\texttt{~with~}\langle v,h\rangle
\mapsto s)$ la preuve de $Q$ obtenue par élimination du $\exists$.

{\footnotesize (\textbf{N.B.} $v$ peut apparaître dans $s$ mais pas
  dans $\Gamma$ ni $Q$.)}
\[
\inferrule{\Gamma\vdash z \;:\; \exists (v:I).\, P\\\Gamma,\, v:I,\, h:P\vdash s \;:\; Q}{\Gamma\vdash (\texttt{match~}z\texttt{~with~}\langle v,h\rangle \mapsto s) \;:\; Q}
\]

\medskip

\itempoint Ceci est conforme à l'idée de BHK : une preuve de
$\exists(v:I).\, P(v)$ est la donnée d'un $t$ de type $I$ d'une preuve
de $Q(t)$.

\end{frame}
%
\begin{frame}
\label{natural-deduction-rules-with-quantifiers-and-notations}
\frametitle{Récapitulatif des notations}

{\footnotesize(À comparer au
  transp. \ref{natural-deduction-rules-with-quantifiers}.)\par}

\medskip

\begin{tabular}{c|c|c}
&Intro&Élim\\\hline
$\Rightarrow$
&\scalebox{0.65}{$\inferrule{\Gamma,v:P\vdash s:Q}{\Gamma\vdash \lambda(v:P).\,s : P\Rightarrow Q}$}
&\scalebox{0.65}{$\inferrule{\Gamma\vdash f:P\Rightarrow Q\\\Gamma\vdash z:P}{\Gamma\vdash fz:Q}$}
\\\hline
$\land$
&\scalebox{0.65}{$\inferrule{\Gamma\vdash z_1:Q_1\\\Gamma\vdash z_2:Q_2}{\Gamma\vdash \langle z_1,z_2\rangle : Q_1\land Q_2}$}
&\scalebox{0.65}{$\inferrule{\Gamma\vdash z:Q_1\land Q_2}{\Gamma\vdash \pi_1 z:Q_1}$}
\quad\scalebox{0.65}{$\inferrule{\Gamma\vdash z:Q_1\land Q_2}{\Gamma\vdash \pi_2 z:Q_2}$}
\\\hline
$\lor$
&\scalebox{0.65}{$\inferrule{\Gamma\vdash z:Q_i}{\Gamma\vdash \iota_i^{(Q_1,Q_2)}z : Q_1\lor Q_2}$} {\footnotesize ($i\in\{1,2\}$)}
&\scalebox{0.65}{$\inferrule{\Gamma\vdash r:P_1\lor P_2\\\Gamma,h_1:P_1\vdash s_1:Q\\\Gamma,h_2:P_2\vdash s_2:Q}{\Gamma\vdash (\texttt{match~}r\texttt{~with~}\iota_1 h_1 \mapsto s_1,\; \iota_2 h_2 \mapsto s_2):Q}$}
\\\hline
$\top$
&\scalebox{0.65}{$\inferrule{\strut}{\Gamma\vdash \bullet:\top}$}
&\scalebox{0.65}{(néant)}
\\\hline
$\bot$
&\scalebox{0.65}{(néant)}
&\scalebox{0.65}{$\inferrule{\Gamma\vdash r:\bot}{\Gamma\vdash \texttt{exfalso}^{(Q)} r:Q}$}
\\\hline
$\forall$
&\scalebox{0.80}{$\inferrule{\Gamma, v:I\vdash s:Q}{\Gamma\vdash \lambda(v:I).\,s\;:\;\forall (v:I).\, Q}$}
&\scalebox{0.80}{$\inferrule{\Gamma\vdash f \;:\; \forall (v:I).\, Q\\{\color{gray}\Gamma\vdash t:I}}{\Gamma\vdash ft \;:\; Q[v\backslash t]}$}
\\\hline
$\exists$
&\scalebox{0.80}{$\inferrule{{\color{gray}\Gamma\vdash t:I}\\\Gamma\vdash z \;:\; Q[v\backslash t]}{\Gamma\vdash \langle t,z\rangle \;:\; \exists (v:I).\, Q}$}
&\scalebox{0.80}{$\inferrule{\Gamma\vdash z \;:\; \exists (v:I).\, P\\\Gamma,\, v:I,\, h:P\vdash s \;:\; Q}{\Gamma\vdash (\texttt{match~}z\texttt{~with~}\langle v,h\rangle \mapsto s) \;:\; Q}$}
\\
\end{tabular}

\medskip

\itempoint Les séquents en gris sont des formations de termes
d'individus (cf. transp. suivant).

\end{frame}
%
\begin{frame}
\frametitle{Monde des termes et monde logique}

\itempoint Les règles pour les démonstrations écrites notamment
transp. \ref{natural-deduction-rules-with-quantifiers} à \ref{natural-deduction-rules-with-quantifiers-and-notations}
\alert{ne sont pas complètes}, il manque les explications sur les
séquents en gris (formation des termes).

\medskip

\itempoint Selon le système logique, on peut distinguer deux
« mondes » \alert{plus ou moins séparés ou confondus} :
\begin{itemize}
\item le monde des termes (individus) et types (d'individus),
\item le monde logique, avec preuves et propositions.
\end{itemize}

\medskip

\itempoint Les règles données aux transp. précédents sont les règles
de construction des \alert{démonstrations} ($\rightarrow$ monde
logique), où se placent tous les séquents sauf ceux marqués en gris.

\medskip

\itempoint Les règles du monde des termes peuvent être calquées sur le
monde des démonstrations, plus simples, ou différentes.

\medskip

\itempoint En Coq/Rocq, les deux mondes sont \alert{séparés mais
  parallèles} : $\mathit{Prop}$ pour le type des propositions et
$\mathit{Type}$ pour le type des types d'individus.

\end{frame}
%
\begin{frame}
\frametitle{Que peut-on quantifier au juste ?}

\itempoint Outre la question des termes d'individus et leur séparation
du monde logique, il manque les explications sur ce qu'on a le droit
de quantifier et d'abstraire ; notamment :
\begin{itemize}
\item peut-on former des propositions comme
  $\forall(A:*).\,(A\Rightarrow A)$ {\footnotesize (où « $*$ » est le
  type des propositions)} avec preuve $\lambda(A:*).\,\lambda(h:A).h$,
  i.e., quantifier sur les propositions (et abstraire dessus dans les
  termes) ?
\item peut-on former des objets comme $\lambda(A:*).\,A$ de type
  $*\rightarrow *$, i.e., abstraire sur les propositions ?
\item peut-on former des propositions comme $\forall(x:I).\,A(x)$ où
  $A$ a pour type $I \rightarrow *$, i.e., quantifier et abstraire sur
  des individus ?
\end{itemize}

\medskip

\itempoint Différents systèmes logiques diffèrent dans la réponse à
ces questions, notamment, les $8$ systèmes du « $\lambda$-cube » de
Barendregt (jusqu'à mélanger complètement preuves et individus).

\end{frame}
%
\begin{frame}
\frametitle{Le problème du $\exists$ et des types sommes}

Doit-on croire à ceci (pour $U$ et $V$ deux types) ?
\[
(\forall(x:U).\,\exists(y:V).\,P(x,y))
\Rightarrow
(\exists(f:U\Rightarrow V).\,\forall(x:U).\,P(x,f(x)))
\]
{\footnotesize
\[
\begin{aligned}
\text{« preuve(?) » :~}
\lambda(h:\cdots).\,
\langle
&\lambda(x:U).\,
(\texttt{match~}hx\texttt{~with~}\langle v,z\rangle \mapsto v),\\
&\lambda(x:U).\,
(\texttt{match~}hx\texttt{~with~}\langle v,z\rangle \mapsto z)
\rangle
\end{aligned}
\]
\par}

\medskip

{\footnotesize (Cet énoncé porte le nom d'\alert{axiome du choix} :
  c'est un analogue pour la théorie des types de l'axiome du choix (de
  Zermelo) en théorie des ensembles.)\par}

\medskip

\itempoint Si on voit $\forall$ et $\exists$ comme des types produit
et \alert{somme} en famille respectivement, \alert{oui} :
$\forall(x:U).\,\exists(y:V).\,P(x,y)$ représente une fonction qui
prend un $x$ de type $U$ et renvoie un $y$ de type $V$ ainsi qu'un
$P(x,y)$ correspondant : on peut collecter tous ces $y$ en une
fonction $f : U \Rightarrow V$.

\medskip

\itempoint Si on voit $\exists$ comme un quantificateur
\alert{logique}, alors \alert{non} : le $y$ renvoyé par $\exists$ ne
peut servir qu'à l'intérieur d'une preuve, pas être collecté en une
fonction.

\medskip

{\footnotesize

\itempoint C'est ici la différence principale entre des systèmes comme
Coq (où l'énoncé ci-dessus ne sera pas prouvable pour $P : U\times V
\to \mathit{Prop}$) et les systèmes à la Martin-Löf comme Agda (où
Curry-Howard est suivi « jusqu'au bout » : il n'y a pas de $\exists$
uniquement logique, et cet énoncé est prouvable comme indiqué).

\par}

\end{frame}
%
\begin{frame}
\frametitle{Imprédicativité}

\itempoint On appelle \textbf{imprédicativité} la possibilité de
définir une proposition ou un type en quantifiant sur toutes les
propositions ou types \alert{y compris celui qu'on définit} : c'est
une forme de circularité.

\medskip

\itempoint P.ex., $\forall (Z:*).\, (Z \Rightarrow A)$ représente le
type des fonctions capables de renvoyer un type $A$ à partir d'un type
$Z$ quelconque, y compris celui qu'on définit.

\smallskip

Cette imprédicativité est utile pour définir des constructions sur les
types.

\medskip

{\footnotesize

Exemples (informellement, et en notant « $*$ » le « type des types »
imprédicatif) :
\begin{itemize}
\item $A \; \cong \; \forall (Z:*).\, (Z \Rightarrow A)$ : donné une
  valeur $x$ de type $A$ on peut en fabriquer une de type
  $Z\Rightarrow A$ comme $\lambda(z:Z).\, x$ pour tout type $Z$, mais
  réciproquement, donné une valeur de type $\forall (Z:*).\, (Z
  \Rightarrow A)$ on peut l'appliquer à $Z = \top$ pour obtenir une
  valeur de type $A$.
\item $A \; \cong \; \forall (Z:*).\, ((A \Rightarrow Z) \Rightarrow
  Z)$ : dans un sens on fabrique $\lambda(k:A \Rightarrow Z).\, kx$
  comme pour le CPS, dans l'autre sens, appliquer à $Z = A$ et
  l'identité.
\item $\bot \; \cong \; \forall (Z:*).\, Z$
\quad\itempoint $A\land B \; \cong \; \forall (Z:*).\, ((A \Rightarrow B
  \Rightarrow Z) \Rightarrow Z)$
\item $A\lor B \; \cong \; \forall (Z:*).\, ((A \Rightarrow Z)
  \Rightarrow (B\Rightarrow Z) \Rightarrow Z)$
\end{itemize}

\par}

\medskip

\itempoint Cela \alert{peut} donner des incohérences logiques
(paradoxe de Girard).

\end{frame}
%
\section{Logique du premier ordre}
\begin{frame}
\frametitle{Logique du premier ordre : principe}

\itempoint La \textbf{logique du premier ordre} ou \textbf{calcul des
  prédicats} (du 1\textsuperscript{er} ordre) est la plus simple qui
ajoute les quantificateurs.  Les « choses » sur lesquelles on a le
droit de quantifier s'appellent des \textbf{individus}.

\medskip

\itempoint Côté typage, elle n'est pas très heureuse : les
« individus » apparaissent comme un type $I$ unique, \textit{ad hoc},
qu'on ne peut presque pas manipuler (la logique ne permet pas de faire
des couples, fonctions, etc., des individus).

\medskip

\itempoint Néanmoins, elle a une \alert{grande importance
  mathématique} car le dogme « orthodoxe » est que :
\begin{center}
Les mathématiques se font dans la « théorie des ensembles\\ de
Zermelo-Fraenkel en logique du premier ordre » ($\mathsf{ZFC}$).
\end{center}

Le manque d'expressivité de la logique (pas de couples, fonctions,
etc.) est \alert{compensé par la théorie elle-même} (constructions
ensemblistes des couples, fonctions, etc.).

\medskip

{\footnotesize\itempoint La \alert{sémantique} (Tarskienne) de la
  logique du premier ordre a aussi des propriétés agréables (théorème
  de complétude de Gödel).\par}

\end{frame}
%
\begin{frame}
\frametitle{Logique du premier ordre : sortes de variables et syntaxe}

\itempoint En (pure) logique du premier ordre, on a diverses sortes de
variables :
\begin{itemize}
\item les \textbf{variables d'individus} ($x$, $y$, $z$...) en nombre
  illimité,
\item les \textbf{variables de prédicats} $n$-aires, ou de
  \textbf{relations} $n$-aires [entre individus] ($A^{(n)}$,
  $B^{(n)}$, $C^{(n)}$...), pour chaque entier naturel $n$.
\end{itemize}

\medskip

\itempoint L'indication d'arité des variables de prédicats est
généralement omise (elle peut se lire sur la formule).

\medskip

\itempoint Une \textbf{formule} (logique) est (inductivement) :
\begin{itemize}
\item l'application $A^{(n)}(x_1,\ldots,x_n)$ d'une variable
  de prédicat à $n$ variables d'individus,
\item l'application d'un connecteur : $(P\Rightarrow Q)$, $(P\land
  Q)$, $(P\lor Q)$ où $P,Q$ sont deux formules, ou encore $\top$,
  $\bot$,
\item une quantification : $\forall x.P$ ou $\exists x.P$
  {\footnotesize (pour $\forall (x:I).P$ ou $\exists (x:I).P$)}, qui
  \alert{lie} la variable d'individu $x$ dans $P$.
\end{itemize}

\medskip

\itempoint\alert{On ne peut quantifier que sur les individus
  (« premier ordre »).}

\end{frame}
%
\begin{frame}
\frametitle{Exemples de formules du premier ordre}

\itempoint Les \textbf{formules propositionnelles} sont encore des
formules du premier ordre, en interprétant chaque variable
propositionnelle comme une variable de prédicat $0$-aire
(« nullaire ») : $A\land B \Rightarrow B\land A$ par exemple.

\bigskip

Autres exemples (qui seront par ailleurs tous démontrables) :
\begin{itemize}
\item $(\forall x.A(x)) \land (\exists x.\top) \Rightarrow (\exists
  x.A(x))$ (ici, $A$ est un prédicat unaire)
\item $(\forall x.\neg A(x)) \Leftrightarrow (\neg\exists x.A(x))$ (idem)
\item $(\exists x.\neg A(x)) \Rightarrow (\neg\forall x.A(x))$ (idem)
\item $(\exists x.A) \Leftrightarrow (\exists x.\top) \land A$ (ici,
  $A$ est un prédicat \alert{nullaire})
\item $(\forall x.A) \Leftrightarrow ((\exists x.\top) \Rightarrow A)$
  (idem)
\item $(\exists x.\forall y.B(x,y))\Rightarrow (\forall y.\exists
  x.B(x,y))$ (ici, $B$ est un prédicat binaire)
\end{itemize}

\bigskip

\textbf{N.B.} On a suivi la convention que $\forall,\exists$ ont une
priorité plus faible que les connecteurs
$\Rightarrow,\lor,\land,\neg$.  Tout le monde n'est pas d'accord avec
cette convention !

\smallskip

\textbf{N.B.2 :} Il serait peut-être préférable de noter $Bxy$ que
$B(x,y)$.

\end{frame}
%
\begin{frame}
\frametitle{Remarques sur la logique du premier ordre}

\itempoint Le type $I$ (non écrit) des « individus », le seul sur
lequel on peut quantifier, est \alert{complètement spécial} en logique
du premier ordre : on ne l'écrit même pas, on ne peut pas former
$I\times I$ ni $I\to I$ ni rien d'autre.

\medskip

\itempoint Ce type $I$ ne se « mélange » pas aux relations
$A,B,C,\ldots$ : les individus vivent dans un monde hermétiquement
séparé des preuves.

\medskip

\itempoint Dans la variante la plus simple, les seuls termes
d'individus sont les variables d'individus (i.e., la seule façon
d'obtenir $t:I$ est d'avoir $t:I$ dans le contexte !).

\medskip

\itempoint On ne suppose pas $I$ habité, i.e. $\exists x.\top$ n'est
pas démontrable (pas plus que $\forall x. A(x) \Rightarrow \exists
x. A(x)$).  Cf. transp. \ref{caveat-inhabited-domain}.

\medskip

\itempoint Pour avoir le droit d'écrire $A(x)$, la variable $x$ doit
avoir été introduite : la règle d'axiome devrait s'écrire
correctement :
\[
\inferrule{ }{\Gamma,\; x:I\vdash x:I}
\quad\text{~et~}\quad
\inferrule{\Gamma\vdash x_1:I\\\cdots\\\Gamma\vdash x_n:I}{\Gamma,\;A(x_1,\ldots,x_n)\vdash A(x_1,\ldots,x_n)}
\]

\end{frame}
%
\begin{frame}
\label{first-order-logic-rules}
\frametitle{Logique du premier ordre : reprise des règles logiques}

\begin{tabular}{c|c|c}
&Intro&Élim\\\hline
$\Rightarrow$
&\scalebox{0.65}{$\inferrule{\Gamma,{\color{mydarkgreen}P}\vdash Q}{\Gamma\vdash P\Rightarrow Q}$}
&\scalebox{0.65}{$\inferrule{\Gamma\vdash P\Rightarrow Q\\\Gamma\vdash P}{\Gamma\vdash Q}$}
\\\hline
$\land$
&\scalebox{0.65}{$\inferrule{\Gamma\vdash Q_1\\\Gamma\vdash Q_2}{\Gamma\vdash Q_1\land Q_2}$}
&\scalebox{0.65}{$\inferrule{\Gamma\vdash Q_1\land Q_2}{\Gamma\vdash Q_1}$}
\quad\scalebox{0.65}{$\inferrule{\Gamma\vdash Q_1\land Q_2}{\Gamma\vdash Q_2}$}
\\\hline
$\lor$
&\scalebox{0.65}{$\inferrule{\Gamma\vdash Q_1}{\Gamma\vdash Q_1\lor Q_2}$}
\quad\scalebox{0.65}{$\inferrule{\Gamma\vdash Q_2}{\Gamma\vdash Q_1\lor Q_2}$}
&\scalebox{0.65}{$\inferrule{\Gamma\vdash P_1\lor P_2\\\Gamma,{\color{mydarkgreen}P_1}\vdash Q\\\Gamma,{\color{mydarkgreen}P_2}\vdash Q}{\Gamma\vdash Q}$}
\\\hline
$\top$
&\scalebox{0.65}{$\inferrule{\strut}{\Gamma\vdash \top}$}
&\scalebox{0.65}{(néant)}
\\\hline
$\bot$
&\scalebox{0.65}{(néant)}
&\scalebox{0.65}{$\inferrule{\Gamma\vdash \bot}{\Gamma\vdash Q}$
(ou pour la logique classique : $\inferrule{\Gamma,\neg Q\vdash \bot}{\Gamma\vdash Q}$)}
\\\hline
$\forall$
&$\inferrule{\Gamma, {\color{mydarkgreen}x:I}\vdash Q}{\Gamma\vdash \forall x.\, Q}$ ($x$ \alert{frais})
&$\inferrule{\Gamma\vdash \forall x.\, Q\\{\color{gray}\Gamma\vdash t:I}}{\Gamma\vdash Q[x\backslash t]}$
\\\hline
$\exists$
&$\inferrule{{\color{gray}\Gamma\vdash t:I}\\\Gamma\vdash Q[x\backslash t]}{\Gamma\vdash \exists x.\, Q}$
&$\inferrule{\Gamma\vdash \exists x.\, P\\\Gamma, {\color{mydarkgreen}x:I}, {\color{mydarkgreen}P}\vdash Q}{\Gamma\vdash Q}$
($x$ \alert{frais})
\\
\end{tabular}

\itempoint Les hypothèses en gris $\Gamma \vdash t:I$ (formation des
termes) ne sont parfois pas écrites dans les arbres de démonstration,
mais sont essentielles (cf. transp. \ref{caveat-inhabited-domain}).

\end{frame}
%
\begin{frame}
\frametitle{Exemple de preuve en logique du premier ordre}

{\footnotesize

\[
\inferrule*[left={$\Rightarrow$Int}]{
\inferrule*[Left={$\forall$Int}]{
\inferrule*[Left={$\exists$Élim}]{
\inferrule*[Left={Ax}]{ }{\exists x.\forall y.B(x,y) \vdash \exists x.\forall y.B(x,y)}
\\
\inferrule*[Right={$\exists$Int}]{
\inferrule*[Right={$\forall$Élim}]{
\inferrule*[Right=Ax]{ }{x:I,\;\forall y.B(x,y),\;y':I \vdash \forall y.B(x,y)}
}{x:I,\;\forall y.B(x,y),\;y':I \vdash B(x,y')}
}{x:I,\;\forall y.B(x,y),\;y':I \vdash \exists x'.B(x',y')}
}{\exists x.\forall y.B(x,y),\;y':I \vdash \exists x'.B(x',y')}
}{\exists x.\forall y.B(x,y) \vdash \forall y'.\exists x'.B(x',y')}
}{\vdash (\exists x.\forall y.B(x,y))\Rightarrow (\forall y'.\exists x'.B(x',y'))}
\]

}

Présentation avec les seules conclusions :

{\footnotesize

\[
\inferrule*[left={$\Rightarrow$Int(\textcolor{mydarkgreen}{$u$})}]{
\inferrule*[Left={$\forall$Int(\textcolor{mydarkgreen}{$y'$})}]{
\inferrule*[Left={$\exists$Élim(\textcolor{mydarkgreen}{$x$},\textcolor{mydarkgreen}{$v$})}]{
\inferrule*[Left={\textcolor{mydarkgreen}{$u$}}]{ }{\exists x.\forall y.B(x,y)}
\\
\inferrule*[Right={$\exists$Int}]{
\inferrule*[Right={$\forall$Élim}]{
\inferrule*[Right={\textcolor{mydarkgreen}{$v$}}]{ }{\forall y.B(x,y)}
}{B(x,y')}
}{\exists x'.B(x',y')}
}{\exists x'.B(x',y')}
}{\forall y'.\exists x'.B(x',y')}
}{(\exists x.\forall y.B(x,y))\Rightarrow (\forall y'.\exists x'.B(x',y'))}
\]

}

\end{frame}
%
\begin{frame}
\frametitle{Exemple de preuve : présentation drapeau}

{\footnotesize
\begin{flagderiv}[example-1st-order-proof]
\assume{mainhyp}{\exists x.\forall y.B(x,y)}{}
\assume{vary}{y'}{}
\assume{exhyp}{x,\;\forall y.B(x,y)}{}
\step{bare}{B(x,y')}{$\forall$Élim sur \ref{exhyp} et $y'$}
\step{exbare}{\exists x'.B(x',y')}{$\exists$Int sur $x$ et \ref{bare}}
\conclude{extrude}{\exists x'.B(x',y')}{$\exists$Elim sur \ref{mainhyp} de \ref{exhyp} dans \ref{exbare}}
\conclude{mainconc}{\forall y'.\exists x'.B(x',y')}{$\forall$Int de \ref{vary} dans \ref{extrude}}
\conclude{}{(\exists x.\forall y.B(x,y))\Rightarrow(\forall y'.\exists x'.B(x',y'))}{$\Rightarrow$Int de \ref{mainhyp} dans \ref{mainconc}}
\end{flagderiv}
\par}

\smallskip

{\footnotesize « Supposons $\exists x.\forall y.B(x,y)$.  Considérons
  un $y'$ arbitraire.  Considérons un $x$ tel que $\forall y.B(x,y)$.
  En particulier, on a $B(x,y')$.  En particulier, $\exists
  x'.B(x',y')$.  Or on pouvait trouver un tel $x$ car $\exists
  x.\forall y.B(x,y)$, donc on a bien la conclusion $\exists
  x'.B(x',y')$.  Le choix de $y'$ étant arbitraire, $\forall
  y'. \exists x'. B(x',y')$.  Finalement, on a prouvé $(\exists
  x.\forall y.B(x,y))\Rightarrow(\forall y'.\exists
  x'.B(x',y'))$. »\par\strut\par}

\end{frame}
%
\begin{frame}
\label{example-predicate-proof-as-lambda-term}
\frametitle{Notation des preuves par des $\lambda$-termes}

On reprend la preuve donnée aux transp. précédents :

\begin{center}
\scalebox{0.8}{\footnotesize
$
\inferrule*[left={$\Rightarrow$Int}]{
\inferrule*[Left={$\forall$Int}]{
\inferrule*[Left={$\exists$Élim}]{
\inferrule*[Left={Ax}]{ }{\exists x.\forall y.B(x,y) \vdash \exists x.\forall y.B(x,y)}
\\
\inferrule*[Right={$\exists$Int}]{
\inferrule*[Right={$\forall$Élim}]{
\inferrule*[Right=Ax]{ }{x:I,\;\forall y.B(x,y),\;y':I \vdash \forall y.B(x,y)}
}{x:I,\;\forall y.B(x,y),\;y':I \vdash B(x,y')}
}{x:I,\;\forall y.B(x,y),\;y':I \vdash \exists x'.B(x',y')}
}{\exists x.\forall y.B(x,y),\;y':I \vdash \exists x'.B(x',y')}
}{\exists x.\forall y.B(x,y) \vdash \forall y'.\exists x'.B(x',y')}
}{\vdash (\exists x.\forall y.B(x,y))\Rightarrow (\forall y'.\exists x'.B(x',y'))}
$
}
\end{center}

Avec les notations du
transp. \ref{natural-deduction-rules-with-quantifiers-and-notations}
elle se note :
\[
\lambda(u:\exists x.\forall y.B(x,y)).\;
\lambda(y':I).\;
(\texttt{match~}u\texttt{~with~}\langle x,v\rangle \mapsto
\langle x,vy'\rangle)
\]

\end{frame}
%
\begin{frame}
\frametitle{Exemples de preuves écrites comme $\lambda$-termes}

\itempoint $(\forall x.(A(x)\Rightarrow C)) \Rightarrow ((\exists
x. A(x)) \Rightarrow C)$

\smallskip

{\footnotesize Preuve :} $\lambda(f:\forall x.(A(x)\Rightarrow C)).\,
\lambda(p:\exists x. A(x)).\, (\texttt{match~}p\texttt{~with~}\langle
x,w\rangle \mapsto fxw)$

\medskip

\itempoint $((\exists x. A(x)) \Rightarrow C) \Rightarrow (\forall
x.(A(x)\Rightarrow C))$

\smallskip

{\footnotesize Preuve :} $\lambda(g:(\exists x. A(x)) \Rightarrow
C).\, \lambda(x:I).\, \lambda(u:A(x)).\, g\langle x,u\rangle$

\medskip

{\footnotesize

\itempoint Notamment pour $C$ valant $\bot$ on a prouvé $(\forall
x.\neg A(x)) \Leftrightarrow \neg(\exists x. A(x))$ ci-dessus.

\par}

\medskip

\itempoint $(\exists x.(A(x)\Rightarrow C)) \Rightarrow ((\forall
x. A(x)) \Rightarrow C)$

\smallskip

{\footnotesize Preuve :} $\lambda(p:\exists x.(A(x)\Rightarrow C)).\,
\lambda(f:\forall x. A(x)).\, (\texttt{match~}\!p\!\texttt{~with~}\!\langle
x,w\rangle \mapsto w(fx))$

\medskip

{\footnotesize

\itempoint Notamment pour $C$ valant $\bot$ on a prouvé $(\exists
x.\neg A(x)) \Rightarrow \neg(\forall x. A(x))$ ci-dessus.

\par}

\medskip

\itempoint $(\forall x.A(x)) \Rightarrow (\exists x.\top) \Rightarrow
(\exists x.A(x))$

{\footnotesize Preuve :} $\lambda(f:\forall x.A(x)).\,
\lambda(p:\exists x.\top).\, (\texttt{match~}p\texttt{~with~}\langle
x,u\rangle \mapsto \langle x,fx\rangle)$

\medskip

\itempoint $\forall z.\exists x.\top$ (cf. transp. suivant)

{\footnotesize Preuve :} $\lambda(z:I).\,\langle z,\bullet\rangle$

\end{frame}
%
\begin{frame}
\label{caveat-inhabited-domain}
\frametitle{Pourquoi des variables d'individus avec les hypothèses ?}

\itempoint L'introduction d'une variable d'individu libre porte en
elle l'hypothèse que \alert{l'univers des individus est habité}
($\exists x. \top$).  Ce fait \alert{n'est pas prouvable} sans cette
hypothèse.  On a $z:I \vdash \exists x.\top$ (ici $z$ variable qcque)
mais \alert{on n'a pas} $\vdash \exists x.\top$.

\medskip

\itempoint Exiger que de pouvoir former $t:I$ à partir de $\Gamma$
permet d'écarter la démonstration \alert{incorrecte} suivante :
\[
\inferrule*[left={$\exists$Int}]{
\inferrule*[Left={$\top$Int}]{ }{\vdash\top}
}{\vdash\exists x.\top}
\]

\medskip

\itempoint En revanche, celle-ci \alert{est correcte} (en utilisant le
terme $z$ pour $t$ dans $\exists$Int) :
\[
\inferrule*[left={$\forall$Int}]{
\inferrule*[Left={$\exists$Int}]{
{\color{gray}\inferrule*[Left={Ax}]{ }{z:I\vdash z:I}}\\
\inferrule*[Right={$\top$Int}]{ }{z:I\vdash\top}
}{z:I\vdash\exists x.\top}
}{\vdash\forall z.\exists x.\top}
\]

\medskip

\textbf{N.B.} Ces problèmes n'ont rien à voir avec la logique
intuitionniste, ils sont identiques en logique classique.  On
\alert{n'a pas} non plus $\forall x.A(x) \Rightarrow \exists x.A(x)$.

\end{frame}
%
\begin{frame}
\frametitle{Monde des individus et monde logique}

\itempoint En logique du premier ordre, on a
\alert{deux « mondes » complètement séparés} :
\begin{itemize}
\item le monde des individus, avec un seul type ($I$) et des variables
  sur lesquelles on peut quantifier,
\item le monde logique, avec propositions et preuves.
\end{itemize}

\medskip

\itempoint Les propositions ont « moralement » un type (qu'on pourrait
appeler « $*$ » ou « $\mathit{Prop}$ »), mais on ne l'écrit pas.  Les
relations $n$-aires ont « moralement » le type $I^n \to
\mathit{Prop}$, pas non plus écrit.

\medskip

\itempoint Dans le transparent \ref{first-order-logic-rules}, tous les
séquents concernent le monde logique (= construction des
démonstrations), sauf ceux marqués en gris.

\medskip

\itempoint Dans la version la plus simple de la logique du premier
ordre (pas de \emph{fonctions} d'individus, seulement des
\emph{relations}), la seule règle du monde des individus est :
\[
\inferrule{ }{\Gamma,\; x:I\vdash x:I}
\]
(on ne peut former un terme d'individu qu'en invoquant une variable du
contexte).

\end{frame}
%
\begin{frame}
\frametitle{Curry-Howard pour la logique du premier ordre}

\itempoint Il faut penser Curry-Howard dans le sens
preuve $\mapsto$ programme.

\smallskip

{\footnotesize (Faute de description précise de règles de typage on ne
  peut pas espérer mieux ici.)\par}

\medskip

\itempoint Curry-Howard va mélanger le monde logique avec le monde des
individus.

\medskip

\itempoint On convertit les propositions en types :
\begin{itemize}
\item $\Rightarrow,\land,\lor,\top,\bot$ deviennent
  $\rightarrow,\times,+,1,0$ comme en calcul propositionnel,
\item $\forall,\exists$ deviennent produits et sommes $\prod,\sum$
  paramétrés par $v:I$.
\end{itemize}

\smallskip

\itempoint On convertit preuves en programmes selon l'interprétation
fonctionnelle des notations données au
transp. \ref{natural-deduction-rules-with-quantifiers-and-notations}.

\medskip

\itempoint P.ex., la preuve de $(\exists x.\forall
y.B(x,y))\Rightarrow (\forall y'.\exists x'.B(x',y'))$ donnée
transp. \ref{example-predicate-proof-as-lambda-term} :
\[
\lambda(u:\exists x.\forall y.B(x,y)).\;
\lambda(y':I).\;
(\texttt{match~}u\texttt{~with~}\langle x,v\rangle \mapsto
\langle x,vy'\rangle)
\]
devient un programme de type
\[
(\sum_{x:I} \prod_{y:I} \, B(x,y)) \rightarrow
(\prod_{y':I} \sum_{x':I} \, B(x',y'))
\]

\end{frame}
%
\begin{frame}
\label{first-order-equality}
\frametitle{L'égalité au premier ordre}

\itempoint En général on veut travailler en logique du premier ordre
\alert{avec égalité}.  C'est-à-dire qu'on introduit une relation
binaire « $=$ » (notée de façon infixe) sujette aux \alert{axiomes}
suivants :
\begin{itemize}
\item réflexivité : $\mathsf{refl} : \forall x.(x=x)$
\item substitution : $\mathsf{subst}^{(\lambda s.P(s))} : \forall
  x.\forall y.((x=y) \Rightarrow P(x) \Rightarrow P(y))$ pour toute
  formule $P(s)$ ayant une variable d'individu libre $s$.
\end{itemize}

\medskip

{\footnotesize \itempoint La logique du premier ordre montre ici ses
  limites : on n'a pas le droit de quantifier sur $P(s)$ ni même
  d'introduire $\lambda (s:I). P(s)$.  Il faut donc comprendre qu'on a
  un « schéma d'axiomes » de substitution, dont chaque
  $\mathsf{subst}^{(\lambda s.P(s))}$ est une instance (et $\lambda
  s.P(s)$ une notation \textit{ad hoc}).\par}

\medskip

\itempoint Exemple de preuve : la symétrie $\forall x.\forall y.((x=y)
\Rightarrow (y=x))$ est prouvée en appliquant la substitution à $P(s)$
valant « $s=x$ », donc par le $\lambda$-terme
\[
\lambda (x:I). \,
\lambda (y:I). \, \lambda (u:(x=y)). \, \mathsf{subst}^{(\lambda
  s.s=x)} \, x \, y \, u \, (\mathsf{refl}\,x)
\]

\medskip

\itempoint Autre exemple : la transitivité $\forall x.\forall
y.\forall z. ((x=y) \Rightarrow (y=z) \Rightarrow (x=z))$ par
\[
\lambda (x:I). \, \lambda (y:I). \, \lambda (z:I). \, \lambda
(u:(x=y)). \, \lambda (v:(y=z)). \, \mathsf{subst}^{(\lambda s.x=s)}
\, y \, z \, v \, u
\]

\end{frame}
%
\section{Arithmétique du premier ordre et théorème de Gödel}
\begin{frame}
\frametitle{L'arithmétique de Heyting et de Peano}

\itempoint L'\textbf{arithmétique du premier ordre} est une (tentative
d')axiomatisation des entiers naturels en logique du premier ordre.
Elle est basée sur les \textbf{axiomes de Peano} (transp. suivant).

\medskip

\itempoint On parle d'\textbf{arithmétique de Heyting} ($\mathsf{HA}$)
en logique intuitionniste, et \textbf{de Peano} ($\mathsf{PA}$) en
logique classique (\alert{mêmes axiomes}, seule la logique change).

\medskip

\itempoint Le cadre de base est la logique du premier ordre
\alert{avec égalité} (cf. transp. \ref{first-order-equality}) et avec
des opérations de formation de termes d'individus $0$ (nullaire), $S$
(unaire) et $+$, $\times$, $\vartriangle$ (binaires) :
\begin{itemize}
\item $0$ est un terme,
\quad\itempoint si $m$ est un terme, $(Sm)$ en est un,
\item si $m$, $n$ sont deux termes, $(m+n)$, $(m\times n)$,
  $(m\vartriangle n)$ en sont.
\end{itemize}

\smallskip

Ils sont censés représenter le successeur ($Sn$ désigne $n+1$), la
somme, le produit et l'exponentiation.  On omet les parenthèses comme
d'habitude.  On peut abréger $1 = S0$ et $2 = S(S0)$, etc.

\end{frame}
%
\begin{frame}
\frametitle{Les axiomes de Peano}

On garde les axiomes de l'égalité :
\begin{itemize}
\item $\forall n.(n=n)$
\item $\forall m.\forall n.((m=n) \Rightarrow P(m) \Rightarrow P(n))$
(schéma de substitution)
\end{itemize}

\medskip

Les \textbf{axiomes de Peano} du premier ordre s'y ajoutent :
\begin{itemize}
\item $\forall n.\neg(Sn=0)$
\item $\forall m.\forall n.((Sm=Sn) \Rightarrow (m=n))$
\item $P(0) \Rightarrow (\forall n.(P(n)\Rightarrow P(Sn)))
\Rightarrow (\forall n.P(n))$ (schéma de \alert{récurrence})
\item $\forall m.(m+0 = m)$
\quad\itempoint $\forall m.\forall n.(m+(Sn) = S(m+n))$
\item $\forall m.(m\times 0 = 0)$
\quad\itempoint $\forall m.\forall n.(m\times (Sn) = m\times n + m)$
\item $\forall m.(m\vartriangle 0 = S0)$
\quad\itempoint $\forall m.\forall n.(m\vartriangle (Sn) = m\vartriangle n \times m)$
\end{itemize}

\medskip

Ci-dessus, $P(s)$ désigne une formule ayant une variable d'individu
libre $s$ : la substitution de l'égalité et la récurrence sont des
\alert{schémas d'axiomes} (un axiome pour chaque $P$ possible) car on
ne peut pas quantifier sur $P$ au premier ordre.

\end{frame}
%
\begin{frame}
\label{example-proof-in-heyting-arithmetic}
\frametitle{Exemple de preuve en arithmétique de Heyting}

Montrons que $\forall n. (n=0 \lor \neg n=0)$ (classiquement c'est une
évidence logique, mais c'est aussi démontrable intuitionistement) :

\medskip

\itempoint On procède par récurrence.  Notons par $P(k)$ la formule
$k=0 \lor \neg k=0$ :
\begin{itemize}
\item $P(0)$ vaut car $0=0$ vaut (réflexivité de l'égalité).
\item $P(Sn)$ vaut car $\neg(Sn=0)$ (premier axiome de Peano).  En
  particulier, $P(n) \Rightarrow P(Sn)$.
\item Donc, par récurrence, $\forall n. P(n)$, ce qu'on voulait
  prouver.
\end{itemize}

\medskip

\itempoint Le $\lambda$-terme de cette preuve ressemble à quelque
chose comme ceci : $\mathsf{recurr}^{(\lambda k.(k=0 \lor \neg
  k=0))}\penalty0\, (\iota_1^{(0=0, \neg 0=0)}(\mathsf{refl}\,0))
\penalty0\, (\lambda (n:\mathsf{nat}).\penalty0\,
\lambda(h:(n=0)\lor\neg(n=0)).\penalty0\, \iota_2^{(Sn=0, \neg
  Sn=0)}(\mathsf{succnotz}\,n))$.

\smallskip

{\footnotesize Ici, « $\mathsf{nat}$ » a été mis pour le type des
  individus (entiers naturels), « $\mathsf{succnotz}$ » pour l'axiome
  de Peano qui affirme $\forall n.\neg(Sn=0)$, et
  « $\mathsf{recurr}^{(\lambda k.P(k))}$ » pour celui qui affirme
  $P(0) \Rightarrow \penalty100 (\forall n.(P(n)\Rightarrow P(Sn)))
  \Rightarrow \penalty100 (\forall n.P(n))$.\par}

\end{frame}
%
\begin{frame}
\frametitle{Quelques théorèmes de l'arithmétique du premier ordre}

On peut prouver en arithmétique de Heyting (et notamment, de Peano)
que :
\begin{itemize}
\item l'addition est commutative, associative, a $0$ pour élément
  neutre...,
\item la multiplication est commutative, associative, a $1 = S0$ pour
  élément neutre...,
\item les identités habituelles sur l'addition, la multiplication,
  l'exponentiation,
\item les propriétés basiques de $m\leq n$ défini par $\exists
  k.(n=m+k)$,
\item les propriétés basiques du codage de Gödel $\langle m,n\rangle$
  défini par $m + \frac{1}{2}(m+n)(m+n+1)$,
\item les propriétés basiques des suites finies codées par $\dbllangle
  a_0,\ldots,a_{k-1}\dblrangle := \langle a_0, \langle a_1,
  \langle\cdots,\langle a_{k-1},0\rangle+1\cdots\rangle+1\rangle+1$,
\item l'existence et l'unicité de la division euclidienne,
\item des propriétés arithmétique de base : existence d'une infinité
  de nombres premiers, existence et unicité de la DFP, irrationalité
  de $\sqrt{2}$ ($\forall p. \forall q. ((p\times p = 2\times q \times
  q) \Rightarrow q = 0)$), etc.
\end{itemize}

\end{frame}
%
\begin{frame}
\frametitle{Curry-Howard pour l'arithmétique de Heyting}

\itempoint La formule $m=n$ est une relation binaire sur les entiers
naturels : elle doit devenir un \alert{type} (paramétré par $m$, $n$,
et habité seulement lorsqu'ils sont égaux) sous l'effet de
Curry-Howard.

\smallskip

\itempoint Il faut y penser comme le type des \textbf{témoignages
  d'égalité} de $m$ et $n$.  En pratique, ce sera un type ayant
seul habitant $\{m\}$ lorsque $m=n$ et aucun sinon.

\medskip

\itempoint Chaque axiome de Peano doit devenir un
programme {\footnotesize (à penser comme l'API d'une bibliothèque
  « entiers naturels »)}.  Le seul non trivial est le schéma de
récurrence $P(0) \Rightarrow (\forall n.(P(n)\Rightarrow P(Sn)))
\Rightarrow (\forall n.P(n))$ : il faut y penser comme la

\alert{primitive récursion} d'une fonction, qui à $c \in A_0$ et
$f(n,\emdash)\colon A_n \to A_{n+1}$ associe la suite $u_n \in
\prod_{n\in\mathbb{N}} A_n$ définie par $u_0 = c$ et $u_{n+1} =
f(n,u_n)$, ou en OCaml

\smallskip

{\footnotesize\tt
let recurr = fun c -> fun f -> let rec u = fun n -> if n==0 then c
else f (n-1) (u (n-1)) in u ;;\newline
{\color{purple}val recurr : 'a -> (int -> 'a -> 'a) -> int -> 'a = <fun>}
\par}

\smallskip

…mais avec un type qui permet à chaque $u_n$ d'être dans un $A_n$
différent : $A_0 \rightarrow (\prod_n\,(A_n\rightarrow A_{n+1}))
\rightarrow (\prod_n\,A_n)$.

\end{frame}
%
\begin{frame}
\frametitle{Curry-Howard pour l'arithmétique de Heyting (2)}

\textcolor{brown}{À quoi ressemble le programme associé à une preuve
  dans l'arithmétique de Heyting ?}

\medskip

On peut souvent s'en faire une idée d'après son type, p.ex. :
\begin{itemize}
\item La commutativité de la multiplication $\forall m.\forall
  n.(m\times n = n\times m)$ prend $m$ et $n$ et un renvoie un
  témoignage d'égalité de $m\times n$ et $n\times m$ (c'est-à-dire en
  fait $m\times n$ calculé de deux manières différentes).
\item La preuve de $\forall n. (n=0 \lor \neg n=0)$ donnée
  transp. \ref{example-proof-in-heyting-arithmetic} prend en entrée
  $n$ et renvoie un type somme avec soit un témoignage d'égalité de
  $n$ à $0$ soit un programme qui donné un tel témoignage renvoie qqch
  d'impossible.  Donc en pratique, ce programme prend $n$ et teste si
  $n=0$.
\item Une preuve de $\forall m.\exists n. Q(m,n)$ va correspondre à un
  programme qui prend $m$ et renvoie $n$ ainsi qu'un témoignage de
  $Q(m,n)$.
\end{itemize}

\medskip

Notamment, \alert{si} $\forall m.\exists n. Q(m,n)$ est prouvable dans
l'arithmétique de \alert{Heyting}, \alert{alors} on peut en déduire
$\varphi_e$ générale récursive totale telle que $\forall
m. Q(m,\varphi_e(m))$ (\textbf{extraction} de programme à partir de la
preuve).

\end{frame}
%
\begin{frame}
\frametitle{Propriétés métamathématiques de l'arithmétique de Heyting}

\itempoint L'arithmétique de Heyting a la \textbf{propriété de la
  disjonction} : si elle prouve $Q_1\lor Q_2$, alors elle prouve $Q_1$
ou $Q_2$.

\smallskip

\itempoint Et la \textbf{propriété de l'existence} : si elle prouve
$\exists n. Q(n)$, alors elle prouve $Q(n)$ pour un $n$ explicite.

\smallskip

{\tiny Petits caractères : ces faits, comme l'extraction de programme,
  dépendent d'un résultat de normalisation sur l'arithmétique de
  Heyting (donc de $\mathsf{Consis}(\mathsf{HA})$).\par}

\medskip

\centerline{*}

\medskip

\itempoint Soit $P^{\cps}$ la formule obtenue en ajoutant
« $\neg\neg$ » devant la formule tout entière, devant la conclusion de
chaque $\Rightarrow$, et après chaque $\forall k$.  Alors Heyting
prouve $P^{\cps}$ ssi Peano prouve $P$ :
\[
\mathsf{PA} \vdash P\text{~\alert{ssi}~}\mathsf{HA} \vdash P^{\cps}
\]

\end{frame}
%
\begin{frame}
\frametitle{Une différence entre Heyting et Peano}

\itempoint On peut formaliser les fonctions générales récursives (ou
machines de Turing) en arithmétique de Heyting.  Par exemple,
$\varphi_e(i){\downarrow}$ signifie $\exists n. T(n,e, i)$ où $T$ est
le prédicat (p.r.) de la forme normale de Kleene, « $n$ code un arbre
de calcul valable de $\varphi_e$ sur l'entrée $i$ ».

\medskip

%% \itempoint Si l'arithmétique de Heyting $\mathsf{HA}$ prouve $\forall
%% m. (Q_1(m) \lor Q_2(m))$, alors il existe $\varphi_e \colon \mathbb{N}
%% \to \{1,2\}$ générale récursive totale telle que $\mathsf{HA}$ prouve
%% $\forall m. (\varphi_e(m){\downarrow}=1 \land Q_1(m) \; \lor \penalty0\;
%% \varphi_e(m){\downarrow}=2 \land Q_2(m))$.

%% \smallskip

\itempoint Si $\mathsf{PA}$ prouve $\forall
m. \varphi_e(m){\downarrow}$, alors $\mathsf{HA}$ le prouve (la
réciproque est évidente).

\medskip

\itempoint Si $\mathsf{HA}$ prouve $\forall m.\exists n. Q(m,n)$,
alors il existe $e$ telle que $\mathsf{HA}$ prouve $\forall
m. \varphi_e(m){\downarrow}$ et $\forall m. Q(m,\varphi_e(m))$.

\bigskip

\itempoint La formule
\[
\forall e. \forall i. (\varphi_e(i){\downarrow} \lor \neg \varphi_e(i){\downarrow})
\]

{\footnotesize (c'est-à-dire $\forall e. \forall i. ((\exists
  n. T(n,e, i)) \lor \neg (\exists n. T(n,e, i)))$)} est évidemment
démontrable dans l'arithmétique de Peano (= en logique classique).
Elle \alert{n'est pas démontrable} en arithmétique de Heyting (= en
logique intuitioniste), car par Curry-Howard on pourrait extraire de
la preuve un algorithme résolvant le problème de l'arrêt.

\smallskip

{\tiny Petits caractères : ces faits, comme l'extraction de programme,
  dépendent d'un résultat de normalisation sur l'arithmétique de
  Heyting (donc de $\mathsf{Consis}(\mathsf{HA})$).\par}

\end{frame}
%
\begin{frame}
\label{proof-checking-predicate}
\frametitle{Formalisation de la démontrabilité}

\textcolor{blue}{Idée-clé :} tester \alert{si une preuve est valable}
est algorithmiquement \alert{décidable} (même primitif récursif).

\smallskip

En revanche, tester si un \alert{énoncé est un théorème} est seulement
\alert{semi-décidable} (en parcourant toutes les preuves possibles).

\medskip

Plus précisément :

\smallskip

\itempoint On peut construire un codage de Gödel pour les formules
arithmétiques et les preuves dans l'arithmétique de Heyting (ou de
Peano), et notamment écrire un prédicat \alert{primitif récursif}
\[
\mathsf{Pf}_{\mathsf{HA}}(n,k)
\quad\text{~resp.~}\quad
\mathsf{Pf}_{\mathsf{PA}}(n,k)
\]
qui signifie « $n$ est le code de Gödel d'une preuve dans
l'arithmétique de Heyting (resp. Peano) de la formule codée par $k$ ».

\medskip

\itempoint Notamment, $\exists n.\, \mathsf{Pf}(n,k)$ peut se lire
comme « $k$ code un théorème », et l'ensemble de ces $k$ est (au
moins) semi-décidable.

\end{frame}
%
\begin{frame}
\label{halting-implies-halting-provability}
\frametitle{Démontrabilité de l'arrêt}

\textcolor{blue}{Idée-clé :} si une machine de Turing s'arrête, on
peut démontrer (dans l'arithmétique de Heyting) qu'elle s'arrête en
donnant une trace d'exécution pas à pas.

\medskip

Plus précisément :

\smallskip

\itempoint Si $T(n,e,i)$, i.e., si $n$ est un arbre de calcul de
$\varphi_e$ sur l'entrée $i$, alors on peut de façon algorithmique
(même p.r.) tirer de $n$ une preuve de $T(n,e,i)$ dans l'arithmétique
de Heyting (i.e., un $n'$ tel que $\mathsf{Pf}_{\mathsf{HA}}(n',k)$ où
$k$ est le code de Gödel de l'énoncé $\varphi_e(i){\downarrow}$,
i.e. $\exists n. T(n,e, i)$).

\medskip

\begin{center}
\alert{Si une machine de Turing s'arrête, alors le fait qu'elle
  s'arrête est prouvable} (dans l'arithmétique de Heyting, \textit{a
  fortiori} de Peano).
\end{center}

\medskip

{\footnotesize\itempoint On notera aussi que si programme fait une
  boucle infinie \emph{explicite évidente}, alors le fait qu'il ne
  termine pas est également prouvable.\par}

\end{frame}
%
\begin{frame}
\frametitle{Gödel, Rosser et Turing jouent ensemble}

{\footnotesize (Preuve du théorème de Gödel revue et corrigée par
  Turing et par Rosser.)\par}

\medskip

Soit $g$ le programme suivant :
\begin{itemize}
\item $g$ cherche en parallèle une preuve (dans l'arithmétique de
  Peano, disons) de l'énoncé « le programme $g$ termine »
  (i.e. $\varphi_g(0){\downarrow}$) et de l'énoncé « le programme $g$
  ne termine pas »,
\item c'est-à-dire qu'il énumère les entiers et, pour chacun, teste
  s'il est le code de Gödel d'une preuve de $\varphi_g(0){\downarrow}$
  ou de $\neg\varphi_g(0){\downarrow}$,
\item s'il trouve (en premier) une preuve que $g$ termine, alors il
  fait une boucle infinie explicite,
\item s'il trouve (en premier) une preuve que $g$ ne termine pas,
  alors il termine immédiatement.
\end{itemize}

\smallskip

Ce programme $g$ a bien un sens, comme d'habitude, par l'« astuce de
Quine » (théorème de récursion de Kleene) + le fait que la
vérification des preuves est algorithmique
(transp. \ref{proof-checking-predicate}).

\end{frame}
%
\begin{frame}
\frametitle{Le théorème de Gödel}

Admettons provisoirement ce qu'on notera
« $\mathsf{Consis}(\mathsf{PA})$ » :
\[
\text{l'arithmétique de Peano ne prouve pas $\bot$}
\]

\bigskip

\itempoint Si le programme $g$ trouve une preuve qu'il ne termine pas,
alors il termine.  Mais ce point donne une preuve qu'il termine
(transp. \ref{halting-implies-halting-provability}).  Donc on a une
preuve de $\bot$ dans l'arithmétique de Peano, contredisant le point
ci-dessus.

\medskip

\itempoint Si $g$ trouve une preuve qu'il termine, il fait une boucle
infinie.  Mais ce point donne une preuve que $g$ ne termine pas
(boucle infinie explicite).  Donc on a une preuve de $\bot$ dans
l'arithmétique de Peano, contredisant le point ci-dessus.

\medskip

Conclusion : $g$ ne trouve ni de preuve qu'il termine ni de preuve
qu'il ne termine pas.  Donc :
\begin{itemize}
\item $g$ ne termine pas,
\item ce fait-là n'est pas prouvable dans l'arithmétique de Peano,
\item mais on l'a prouvé à l'aide de $\mathsf{Consis}(\mathsf{PA})$,
  donc $\mathsf{Consis}(\mathsf{PA})$ lui-même n'est pas prouvable
  dans Peano (toujours en supposant $\mathsf{Consis}(\mathsf{PA})$).
\end{itemize}

\end{frame}
%
\begin{frame}
\frametitle{Cohérence de Peano}

\itempoint L'énoncé $\mathsf{Consis}(\mathsf{PA})$ peut se lire
ainsi :

\smallskip

« Le programme $g'$ qui parcourt les entiers et, pour chacun, teste
s'il est le code de Gödel d'une preuve de $\bot$ dans l'arithmétique
de Peano et dans ce cas termine, ne termine pas. »

\smallskip

Cet énoncé \alert{a un sens} dans l'arithmétique du premier ordre,
mais (on vient de le voir) \alert{n'est pas démontrable} s'il est
vrai.

\medskip

{\footnotesize

\itempoint L'énoncé $\mathsf{Consis}(\mathsf{HA})$ analogue pour
l'arithmétique de Heyting \alert{est équivalent} à
$\mathsf{Consis}(\mathsf{PA})$ par la traduction CPS (et cette
équivalence est prouvable dans $\mathsf{HA}$).

\smallskip

Notamment, Peano ne prouve pas non plus
$\mathsf{Consis}(\mathsf{HA})$.  Par propriété de la disjonction,
$\mathsf{HA}$ ne prouve même pas $\mathsf{Consis}(\mathsf{HA}) \lor
\neg\mathsf{Consis}(\mathsf{HA})$.

\par}

\bigskip

\itempoint En revanche, $\mathsf{ZFC}$ (le cadre usuel pour faire des
mathématiques) démontre $\mathsf{Consis}(\mathsf{PA})$ : « Les axiomes
de Peano sont vrais dans $\mathbb{N}$ donc leurs conséquences le sont
aussi, et notamment $\bot$ ne peut pas en faire partie.  (Et au
passage, si $\mathsf{PA}$ démontre $\varphi_e(i){\downarrow}$, alors
$\varphi_e(i){\downarrow}$ est vrai.) »

\smallskip

Donc $\mathsf{ZFC}$ est strictement plus fort que $\mathsf{PA}$ (même
pour l'arithmétique).

\end{frame}
%
\begin{frame}
\frametitle{Le théorème de Gödel généralisé}

Pour n'importe quelle sorte de « théorie logique » (classique ou
intuitioniste, pas limitée au premier ordre) $T$ telle que :
\begin{itemize}
\item les énoncés et démonstrations sont codables par des entiers
  naturels,
\item on peut algorithmiquement tester si un entier naturel code une
  démonstration valable dans $T$, et quelle est sa conclusion,
\item $T$ permet de formaliser « $\varphi_e(i){\downarrow}$ »
  (calculablement en $e$ et $i$),
\item si $\varphi_e(i){\downarrow}$ alors on peut tirer d'une trace
  d'exécution une preuve de ce fait dans $T$, et idem pour une boule
  infinie explicite,
\end{itemize}
on peut construire le programme $g_T$ qui cherche en parallèle dans
$T$ une preuve de que $g_T$ termine ou ne termine pas, et fait une
boule infinie explicite dans le premier cas, termine immédiatement
dans le second.

\medskip

\itempoint Si $T$ ne prouve pas $\bot$ (hypothèse notée
« $\mathsf{Consis}(T)$ »), alors $g_T$ ne termine pas, mais $T$ ne
peut pas le prouver.  Notamment, $T$ ne prouve pas
$\mathsf{Consis}(T)$ (toujours si $\mathsf{Consis}(T)$).
{\footnotesize Ceci s'applique notamment à Coq, à $\mathsf{ZFC}$,
  etc.}

\end{frame}
%
\begin{frame}
\frametitle{L'ensemble des théorèmes est semi-décidable non décidable}

Avec $T$ comme au transparent précédent (et sous l'hypothèse
$\mathsf{Consis}(T)$), par exemple $\mathsf{PA}$ supposons par
l'absurde qu'on puisse décider algorithmiquement si une formule $P$
est un théorème de $T$.

\medskip

Soit $g''$ le programme qui :
\begin{itemize}
\item teste si $\varphi_{g''}(0){\downarrow}=0$ est un théorème de $T$
  (grâce à l'hypothèse effectuée),
\item si oui, termine et renvoie $1$,
\item si non, termine et renvoie $0$.
\end{itemize}

\medskip

\itempoint Par construction, $g''$ termine forcément et renvoie soit
$0$ soit $1$.  Si $g''$ termine et renvoie $0$, alors $T$ le prouve,
donc $g''$ termine et renvoie $1$, contradiction ; si $g''$ termine et
renvoie $1$, alors $T$ le prouve, donc (par $\mathsf{Consis}(T)$) il
ne proue pas que $g''$ termine et renvoie $0$, donc $g''$ termine et
renvoie $0$, contradiction.

\end{frame}
%
%% \mode<article>{
%%     \ifodd\value{page}
%%         \clearpage\hbox{}\newpage
%%         \thispagestyle{empty}
%%     \fi}
%
\end{document}
